% Basic LaTeX examples adapted from Leslie Lamport.
\chapter{Introduction}
\label{introduction}


This template demonstrates how to present mathematical statements, equations,
lists, tables, figures, algorithms, code, output, and references in a thesis.
For writing conventions, follow \citet{schifano2026statistical}, available at
\url{https://statds.github.io/stat-writing/}. The University of Connecticut's
requirements govern thesis formatting.


\section{Basic LaTeX examples}
These source examples illustrate conventions described by
\citet{lamport1994latex}.


\subsection{Ordinary text}


The ends of words and sentences are marked by spaces. It doesn't matter how
many spaces you type; one is as good as 100. The end of a line counts as a
space.\footnote{ This is a sample input file. Comparing it with the output it
generates can show you how to produce a simple document of your own. }


One or more blank lines denote the end of a paragraph.


Since any number of consecutive spaces are treated like a single one, the
formatting of the input file makes no difference to \LaTeX, but it makes a
difference to you. When you use \LaTeX, making your input file as easy to read
as possible will be a great help as you write your document and when you change
it. Keep explanatory comments brief and separate them from the thesis text.


Because printing is different from typewriting, there are a number of things
that you have to do differently when preparing an input file than if you were
just typing the document directly. Quotation marks like ``this'' have to be
handled specially, as do quotes within quotes: ``\,`this' is what I just wrote,
not `that'\,.''


Dashes come in three sizes: an intra-word dash, a medium dash for number ranges
like 1--2, and a punctuation dash---like this.


A sentence-ending space should be larger than the space between words within a
sentence. You sometimes have to type special commands in conjunction with
punctuation characters to get this right, as in the following sentence. Gnats,
gnus, etc.\ all begin with G\@. You should check the spaces after periods when
reading your output to make sure you haven't forgotten any special cases.
Generating an ellipsis \ldots\ with the right spacing around the periods
requires a special command.


\LaTeX\ interprets some common characters as commands, so you must type special
commands to generate them. These characters include the following: \$ \& \% \#
\{ and \}.


In printing, text is emphasized by using an {\em italic\/} type style.


\begin{em}
  A long segment of text can also be emphasized in this way. Text within such a
  segment given additional emphasis with\/ {\em Roman} type. Italic type loses
  its ability to emphasize and become simply distracting when used excessively.
\end{em}


It is sometimes necessary to prevent \LaTeX\ from breaking a line where it
might otherwise do so. This may be at a space, as between the ``Mr.'' and
``Jones'' in ``Mr.~Jones,'' or within a word---especially when the word is a
symbol like \mbox{\em itemnum\/} that makes little sense when hyphenated across
lines.


Footnotes\footnote{This is an example of a footnote.} pose no problem.


\LaTeX\ is good at typesetting mathematical formulas like
$\widehat{\mu}=\bar{x}$ or
\[
  \operatorname{SE}(\overline{X}) = \sigma/\sqrt{n}.
\]


Remember that a letter like $x$ is a formula when it denotes a mathematical
symbol, and should be treated as one.


\subsection{Displayed text}


Text is displayed by indenting it from the left margin.


\subsubsection{Quotations}


Quotations are commonly displayed. There are short quotations
\begin{quote}
  This is a short quotation. It consists of a single paragraph of text. There
  is no paragraph indentation.
\end{quote}
and longer ones.
\begin{quotation}
  This is a longer quotation. It consists of two paragraphs of text. The
  beginning of each paragraph is indicated by an extra indentation.


  This is the second paragraph of the quotation. It is just as dull as the
  first paragraph.
\end{quotation}


\subsubsection{Lists}


Another frequently-displayed structure is a list.


\paragraph{Itemize.}
The following is an example of an {\em itemized} list.


\begin{quote}
  \begin{itemize}
    \item  This is the first item of an itemized list.  Each item
    in the list is marked with a ``tick''. The document style determines what
    kind of tick mark is used.


    \item  This is the second item of the list.  It contains another
    list nested inside it. The inner list is an {\em enumerated} list.
    \begin{enumerate}
      \item This is the first item of an enumerated list that
      is nested within the itemized list.


      \item This is the second item of the inner list.  \LaTeX\
      allows you to nest lists deeper than you really should.
    \end{enumerate}
    This is the rest of the second item of the outer list. It is no more
    interesting than any other part of the item.
    \item  This is the third item of the list.
  \end{itemize}
\end{quote}


\paragraph{Verse.}
You can even display poetry.


\begin{quote}
  \begin{quote}
    \begin{verse}
      There is an environment for verse \\
      Whose features some poets will curse.


      For instead of making\\
      Them do {\em all\/} line breaking, \\
      It allows them to put too many words on a line when they'd
      rather be forced to be terse.
    \end{verse}
  \end{quote}
\end{quote}


\subsubsection{Mathematics}
Displayed formulas may occupy one or several lines. Use the environments
provided by \texttt{amsmath} for multiline equations.
\[
  Y_i = \beta_0+\beta_1 x_i+\varepsilon_i.
\]
Don't start a paragraph with a displayed equation, nor make one a paragraph by
itself.


\section{Sampling and estimation}
Suppose a study records independent observations $X_1,\ldots,X_n$ from a
population with mean $\mu$ and variance $\sigma^2$, where $n\geq 1$ and
$0<\sigma^2<\infty$. The estimand is the population mean. An estimator is a
rule applied to the random sample; an estimate is its value for the observed
data. This distinction separates uncertainty before sampling from a number
reported after analysis \citep{wasserman2004all}.


\begin{definition}[Bias]
  The bias of an estimator $T$ of $\mu$, when its expectation exists, is
  \[
    \operatorname{Bias}_{\mu}(T) = \mathbb{E}_{\mu}(T)-\mu.
  \]
  An estimator is unbiased when its bias is zero for every population in the
  specified model.
\end{definition}
The sample mean is the estimator
\begin{equation}
  \label{eq:random-sample-mean}
  \overline{X} = \frac{1}{n}\sum_{i=1}^{n}X_i.
\end{equation}
Equation~\eqref{eq:random-sample-mean} uses uppercase letters for random
variables. Lowercase $x_i$ and $\bar{x}$ denote their observed values.


\begin{theorem}[Properties of the sample mean]
  \label{thm:sample-mean-properties}
  Under the stated sampling assumptions, $\overline{X}$ is unbiased for $\mu$
  and has variance $\sigma^2/n$.
\end{theorem}
\begin{proof}
  Linearity of expectation gives
  \[
    \mathbb{E}(\overline{X})
    = \frac{1}{n}\sum_{i=1}^{n}\mathbb{E}(X_i) = \mu.
  \]
  Independence makes the pairwise covariances zero, so
  \[
    \operatorname{Var}(\overline{X})
    = \frac{1}{n^2}\sum_{i=1}^{n}\operatorname{Var}(X_i)
    = \frac{\sigma^2}{n}.
  \]
\end{proof}
\begin{corollary}[Precision and sample size]
  Under the assumptions of Theorem~\ref{thm:sample-mean-properties}, the
  standard error of $\overline{X}$ is $\sigma/\sqrt{n}$. Multiplying the
  sample size by four halves this standard error.
\end{corollary}
\begin{proof}
  Take the square root of the variance and replace $n$ by $4n$.
\end{proof}
The standard error describes variation across samples. It differs from the
population standard deviation, which describes variation among individuals.
When $\sigma$ is unknown, estimate it rather than reporting the formula as
though the population variance were known.


\section{Assumptions and interpretation}
Unbiasedness alone does not establish that an estimator is useful. Report its
variability, the assumptions behind its uncertainty calculation, and the
population to which the conclusions apply. Independence deserves particular
attention when observations come from repeated measurements or clusters.
Chapter~\ref{chap:normal-mean} develops a confidence interval under an explicit
normal sampling model.


\begin{remark}
  Increasing sample size reduces sampling variability under this model. It
  does not automatically remove selection bias, measurement error, or model
  misspecification.
\end{remark}
